% ==========================================================
%  EXAM QUESTIONS — AI IN LIFE SCIENCES  (WORKED SOLUTIONS)
%  Companion to exams.tex: every question is followed
%  immediately by a green answer box. Correct options are
%  marked with a green check, incorrect ones with a red cross,
%  each with a short justification (incl. why tricky near-miss
%  options are wrong).
%
%  NOTE: exams.tex stays the blank practice set. This file is
%  the solutions variant — do not merge the two.
% ==========================================================

\documentclass[11pt,oneside]{article}

\usepackage[
  a4paper,
  top=2.4cm,
  bottom=2.1cm,
  left=2.3cm,
  right=2.3cm,
  headheight=15pt
]{geometry}

\usepackage[T1]{fontenc}
\usepackage[utf8]{inputenc}
\usepackage{lmodern}
\usepackage{microtype}
\usepackage{inconsolata}
\usepackage{amsmath,amssymb,mathtools,bm}
\usepackage{enumitem}
\usepackage{booktabs,array,multicol,longtable}
\usepackage{xcolor}
\usepackage[most]{tcolorbox}
\usepackage{titlesec}
\usepackage{fancyhdr}
\usepackage[hidelinks]{hyperref}

\setlength{\parindent}{0pt}
\setlength{\parskip}{0.55em}
\setlist[itemize]{topsep=3pt,itemsep=2pt,leftmargin=2.2em}
\setlist[enumerate]{topsep=3pt,itemsep=2pt,leftmargin=2.2em}

% ── Colors ───────────────────────────────────────────────
% Shared accent with summary.tex — do not change.
\definecolor{accent}{HTML}{1B3A5C}
\definecolor{q@frame}{HTML}{1A5276}
\definecolor{q@bg}{HTML}{F6FAFD}
\definecolor{key@frame}{HTML}{1E8449}
\definecolor{key@bg}{HTML}{F4FBF6}
\definecolor{note@frame}{HTML}{9A6A00}
\definecolor{note@bg}{HTML}{FEF9E7}
\definecolor{ok}{HTML}{1E8449}    % green check — correct option
\definecolor{bad}{HTML}{B03A2E}   % red cross  — incorrect option

% ── Box global style ─────────────────────────────────────
\tcbset{
  enhanced,
  breakable,
  arc=2.5mm,
  boxrule=0.5pt,
  left=4.5mm,
  right=4mm,
  top=2.5mm,
  bottom=2.5mm,
  fonttitle=\bfseries\small,
  sharp corners=west,
  parbox=false
}

% Question box (blue) — one per exam question.
\newtcolorbox{qbox}[1][]{
  colback=q@bg,
  colframe=q@frame,
  borderline west={3.5pt}{0pt}{q@frame},
  title={#1}
}

% Key insight box (green) — answer hints, key results.
\newtcolorbox{keybox}[1][]{
  colback=key@bg,
  colframe=key@frame,
  borderline west={3.5pt}{0pt}{key@frame},
  title={#1}
}

% Organizational note box (amber) — structural remarks only.
\newtcolorbox{notebox}[1][]{
  colback=note@bg,
  colframe=note@frame,
  borderline west={3.5pt}{0pt}{note@frame},
  title={#1}
}

% ── Section headings ─────────────────────────────────────
\titleformat{\section}
  {\normalfont\Large\bfseries\color{accent}}
  {}{0em}{}
  [{\vspace{1pt}\color{accent!45}\titlerule[0.45pt]}]
\titlespacing*{\section}{0pt}{3.2ex plus 1ex minus .2ex}{1.8ex plus .2ex}

\titleformat{\subsection}
  {\normalfont\normalsize\bfseries\color{accent!80!black}}
  {}{0em}{}
\titlespacing*{\subsection}{0pt}{2.2ex plus 1ex minus .2ex}{1.0ex plus .2ex}

% ── Header & footer ──────────────────────────────────────
\pagestyle{fancy}
\fancyhf{}
\fancyhead[L]{\small\color{gray!70}Old Exam Questions (Solutions) — \coursetitle}
\fancyhead[R]{\small\color{gray!70}\thepage}
\renewcommand{\headrulewidth}{0.4pt}

% ── Document metadata ─────────────────────────────────────
\newcommand{\coursetitle}{AI in Life Sciences}

% ── Question counter & macros ────────────────────────────
\newcounter{question}
\newcommand{\questiontitle}[1]{%
  \stepcounter{question}\begin{qbox}[Question \thequestion: #1]}
\newcommand{\closequestion}{\end{qbox}}

% Mark questions that appeared on multiple exams.
\newcommand{\repeatnote}{\textit{\color{gray!70}Asked multiple times.}}

% Multiple-choice checkbox bullet.
\newcommand{\choice}{\item[$\square$] }

% ── Inline answer macros (solutions variant) ─────────────
% Green answer box placed directly under each question. The argument records the
% correct letters (kept for reference / grep), but the header only shows "Answers"
% so the key is not spoiled at a glance.
% Usage:
%   \begin{answer}{A, B}
%     \yes{A} reason why correct ...
%     \no{C}  reason why wrong (esp. for tricky near-misses) ...
%   \end{answer}
\newenvironment{answer}[1]{%
  \begin{keybox}[Answers]%
  \begin{itemize}[leftmargin=1.7em,itemsep=1.5pt,topsep=2pt]}%
  {\end{itemize}\end{keybox}}
\newcommand{\yes}[1]{\item[\textcolor{ok}{\ensuremath{\checkmark}}]\textbf{#1.}\ }
\newcommand{\no}[1]{\item[\textcolor{bad}{\ensuremath{\bm\times}}]\textbf{#1.}\ }

% ── Math macros (minimal subset for exam questions) ──────
\newcommand{\R}{\mathbb{R}}
\newcommand{\E}{\mathbb{E}}
\newcommand{\Prob}{\mathbb{P}}
\newcommand{\KL}[2]{D_{\mathrm{KL}}\!\left(#1\,\|\,#2\right)}
\newcommand{\dd}{\,\mathrm{d}}


% ==========================================================
%  BEGIN DOCUMENT
% ==========================================================
\begin{document}

% ── Title page ───────────────────────────────────────────
\thispagestyle{empty}
\vspace*{1.4cm}
\begin{center}
  {\color{accent}\rule{\linewidth}{1.8pt}}\par
  \vspace{0.65cm}
  {\fontsize{26}{32}\selectfont\bfseries\color{accent}Old Exam Questions\\[2pt] \coursetitle\par}
  \vspace{0.25cm}
  {\Large\color{key@frame}\bfseries Worked Solutions\par}
  \vspace{0.55cm}
  {\color{accent!35}\rule{\linewidth}{0.5pt}}\par
  \vspace{0.55cm}
  {\large Stefan Eder\par}
  \vspace{0.15cm}
  {\small\color{gray!70}\today\par}
\end{center}

\vspace{0.9cm}
\begin{notebox}[How this file is organized]
This is the \textbf{solutions companion} to \texttt{exams.tex}. Every question is
followed \emph{immediately} by a green \textbf{answer box}. Inside it, each option is
marked \textcolor{ok}{\ensuremath{\checkmark}} (correct) or
\textcolor{bad}{\ensuremath{\bm\times}} (incorrect), with a short justification ---
\emph{why} the correct options are correct and, for the tricky / near-miss
distractors, \emph{why} they are wrong. Questions come from two sittings: the
\textbf{2024} exam (25 June 2024, Sections~1--6, \emph{multiple-select} --- one
\emph{or more} options correct) and the \textbf{2026} exam (25 June 2026,
Sections~7--10, \emph{single-best-answer} --- exactly one option correct). Wording
and the answer options have both been paraphrased (the meaning and the correct
answers are preserved). For a clean practice set with blank answers, use
\texttt{exams.tex}.
\end{notebox}

\tableofcontents
\newpage


% ==========================================================
%  SECTION 1 — Drug Discovery, Molecules & Representations
% ==========================================================
\section{Drug Discovery, Molecules \& Representations}

\questiontitle{SMILES strings}
Regarding the SMILES notation for molecules, which of the following statements hold?
\begin{itemize}[leftmargin=1.8em]
  \choice SMILES is short for the \emph{Simplified Molecular Input Line Entry System}.
  \choice A SMILES encodes a molecule as a one-dimensional sequence of characters.
  \choice A SMILES string cannot be turned into a molecular graph.
  \choice Each molecule has exactly one SMILES, so no two distinct molecules can share a string.
\end{itemize}
\closequestion
\begin{answer}{A, B}
  \yes{A} It is exactly the \emph{Simplified Molecular Input Line Entry System}.
  \yes{B} A molecule is encoded as a 1D string of characters (atoms, bonds, ring/branch tokens).
  \no{C} SMILES \emph{can} be parsed into a molecular graph --- that round-trip is one of its main uses.
  \no{D} Not unique: a single molecule has \emph{many} valid SMILES. Only a \emph{canonical} SMILES (a fixed algorithm) gives one representative string --- the plain format is not inherently unique.
\end{answer}

\questiontitle{Reading a SMILES string}
Consider the molecule \emph{pyruvic acid} (2-oxopropanoic acid), i.e.\ a methyl
group followed by a ketone carbonyl and then a carboxylic acid
($\text{CH}_3\text{--C}(\text{=O})\text{--C}(\text{=O})\text{--OH}$). Which of the
SMILES strings below correctly encode this molecule?
\begin{itemize}[leftmargin=1.8em]
  \choice \texttt{O=C(C(=O)O)C}
  \choice \texttt{CC(=O)C(=O)O}
  \choice \texttt{CCCOOO}
  \choice \texttt{OCCOO}
\end{itemize}
\closequestion
\begin{answer}{A, B}
  \yes{A} \texttt{O=C(C(=O)O)C}: a carbonyl carbon bearing a \texttt{C(=O)O} (carboxyl) and a \texttt{C} (methyl) --- exactly CH\textsubscript{3}--CO--COOH, just written starting from the ketone.
  \yes{B} \texttt{CC(=O)C(=O)O}: methyl, ketone carbonyl, carboxyl read left-to-right --- the same molecule.
  \no{C} \texttt{CCCOOO} is a chain of three carbons then three oxygens (no double bonds, no carboxyl) --- not pyruvic acid, and chemically nonsensical.
  \no{D} \texttt{OCCOO} encodes a short O/C chain with the wrong connectivity and no ketone/carboxyl --- not the target.
\end{answer}

\questiontitle{Molecular fingerprints}
Which of the following statements about molecular fingerprints are accurate?
\begin{itemize}[leftmargin=1.8em]
  \choice A unique molecular graph can be reconstructed from them.
  \choice They are binary or numeric encodings of molecular structure.
  \choice They serve similarity searches and machine-learning tasks.
  \choice They give a 3D representation of the molecule.
\end{itemize}
\closequestion
\begin{answer}{B, C}
  \no{A} They are \emph{lossy} and often hashed (bit collisions), so you cannot reconstruct a \emph{unique} molecular graph from them.
  \yes{B} Fingerprints are fixed-length \emph{binary} (or count/numerical) vectors encoding the presence of substructures.
  \yes{C} Their classic uses are Tanimoto \emph{similarity search} and as ML features.
  \no{D} They encode 2D substructure/topology, \emph{not} 3D geometry.
\end{answer}

\questiontitle{Bioassays}
Select the correct statements about bioassays.
\begin{itemize}[leftmargin=1.8em]
  \choice They quantify the biological response or function that a substance triggers.
  \choice They usually probe the substance across a range of concentrations.
  \choice They are wet-lab procedures for characterising a compound in a biological setting.
  \choice They are usually carried out \emph{in vitro}, i.e.\ inside a living organism.
\end{itemize}
\closequestion
\begin{answer}{A, B, C}
  \yes{A} The whole point of an assay is to quantify a biological response/function.
  \yes{B} They are run over a concentration range (dose--response), giving e.g.\ IC\textsubscript{50}.
  \yes{C} They are physical wet-lab experiments that generate the bioactivity labels we later model.
  \no{D} The trap is the definition: \emph{in vitro} means ``in glass'' (test tube / cell culture), \emph{not} in a living organism. ``In a living organism'' is \emph{in vivo}.
\end{answer}

\questiontitle{ADME properties}
Which statements about ADME (Absorption, Distribution, Metabolism, and Excretion)
properties are correct?
\begin{itemize}[leftmargin=1.8em]
  \choice Lipinski's rule of five can serve as a proxy for ADME behaviour.
  \choice Favourable ADME properties guarantee that a compound becomes an effective drug.
  \choice Certain ADME properties, such as clearance, are hard to model.
  \choice ADME properties are decisive for a drug's pharmacokinetics.
\end{itemize}
\closequestion
\begin{answer}{A, C, D}
  \yes{A} Lipinski's Rule of Five is a cheap drug-likeness heuristic used as an ADME/oral-bioavailability proxy.
  \no{B} Necessary but \emph{not sufficient}: good ADME means the drug reaches its target intact, but efficacy still needs actual target binding and the right pharmacodynamics.
  \yes{C} Properties like clearance depend on complex metabolism and are notoriously hard to predict.
  \yes{D} ADME \emph{is} pharmacokinetics --- what the body does to the drug over time.
\end{answer}

\questiontitle{Characteristics of QSAR / bioactivity data}
Which of the following correctly characterise typical QSAR / bioactivity datasets?
\begin{itemize}[leftmargin=1.8em]
  \choice QSAR datasets typically contain on the order of thousands to millions of data points.
  \choice QSAR datasets usually show a balanced mix of active and inactive compounds.
  \choice QSAR datasets frequently include groups of structurally related compounds.
  \choice QSAR labels are often noisy because of biological variability and experimental inconsistency.
\end{itemize}
\closequestion
\begin{answer}{C, D}
  \no{A} Bioactivity data is usually \emph{small/limited} (hundreds to a few thousand per target), not millions.
  \no{B} It is typically strongly \emph{imbalanced} --- far more inactives than actives.
  \yes{C} Medicinal-chemistry projects explore \emph{analog series} (one scaffold, small variations), so related compounds cluster heavily.
  \yes{D} Assays vary between labs/days and have measurement error, so labels are noisy.
\end{answer}

\questiontitle{Applications of QSAR modeling}
For which of the following applications is QSAR modeling relevant?
\begin{itemize}[leftmargin=1.8em]
  \choice Building 3D models of protein structures.
  \choice Finding promising drug candidates.
  \choice Estimating the toxicity of novel chemical compounds.
  \choice Mapping the genome sequences of organisms.
\end{itemize}
\closequestion
\begin{answer}{B, C}
  \no{A} 3D protein structure is \emph{structure prediction} (e.g.\ AlphaFold), a different problem --- QSAR maps molecule $\to$ property.
  \yes{B} QSAR ranks/scores compounds by predicted activity --- i.e.\ finds candidates.
  \yes{C} Predicting properties like toxicity (QSPR) is a core QSAR application.
  \no{D} Genome sequencing/mapping is genomics, unrelated to QSAR.
\end{answer}


% ==========================================================
%  SECTION 2 — Virtual Screening & QSAR Modeling
% ==========================================================
\section{Virtual Screening \& QSAR Modeling}

\questiontitle{Ligand-based virtual screening}
Which statements about ligand-based virtual screening are correct?
\begin{itemize}[leftmargin=1.8em]
  \choice Molecular fingerprints are a common ingredient of ligand-based virtual screening.
  \choice Ligand-based virtual screening cannot predict the bioactivity of untested compounds.
  \choice Ligand-based virtual screening needs the 3D structure of the target protein.
  \choice Ligand-based virtual screening selects compounds using bioactivity predictions alone.
\end{itemize}
\closequestion
\begin{answer}{A}
  \yes{A} Fingerprint similarity to known actives is the workhorse of ligand-based VS.
  \no{B} It is precisely \emph{used} to predict activity of unknown/untested compounds by analogy to known actives.
  \no{C} Needing the protein's 3D structure is \emph{structure-based} VS. ``Ligand-based'' means you only use known ligands, no target structure.
  \no{D} Selection also weighs ADME, drug-likeness and structural-alert filters --- not bioactivity alone.
\end{answer}

\questiontitle{Data splits in QSAR model evaluation}
When evaluating QSAR models, which of the following statements about
data-splitting strategies hold?
\begin{itemize}[leftmargin=1.8em]
  \choice Temporal splits gauge performance on future, unseen data by mimicking how data accrue over time.
  \choice A random split is generally the best evaluation scheme because it removes all bias.
  \choice Cross-validation is seldom applied to QSAR because of its computational cost.
  \choice Cluster splits keep similar compounds together so that structurally similar ones do not span the training and test sets.
\end{itemize}
\closequestion
\begin{answer}{A, D}
  \yes{A} Splitting by time mimics how data really arrives and tests genuine prospective generalisation.
  \no{B} A \emph{random} split is the classic trap: structurally similar compounds leak into both sets (compound-series bias), so it \emph{over-estimates} performance rather than eliminating bias.
  \no{C} Cross-validation \emph{is} routinely used in QSAR; cost is not a reason to avoid it here.
  \yes{D} Scaffold/cluster splits keep an entire analog series on one side, so the test set is not just memorised neighbours.
\end{answer}

\questiontitle{Early-recognition metrics in virtual screening}
Because only the top of a virtual-screening ranking is followed up experimentally,
recognising actives early matters most. Which of the following metrics are
specifically designed to reward strong early recognition?
\begin{itemize}[leftmargin=1.8em]
  \choice Precision-Recall Curve.
  \choice BEDROC (Boltzmann-Enhanced Discrimination of ROC).
  \choice Average Precision (AP).
  \choice Enrichment Factor (EF).
\end{itemize}
\closequestion
\begin{answer}{B, D}
  \no{A} A Precision--Recall curve summarises the \emph{whole} ranking, not specifically the top.
  \yes{B} BEDROC applies an exponential weight that emphasises actives found \emph{early} in the ranking.
  \no{C} Average Precision is a general ranking metric and is not early-recognition-weighted.
  \yes{D} The Enrichment Factor measures how many actives appear in the top $x\%$ --- a pure early-recognition metric.
\end{answer}

\questiontitle{Target identification / deconvolution}
Which statements about target identification / deconvolution are correct?
\begin{itemize}[leftmargin=1.8em]
  \choice For a given molecule, it determines which of several simultaneously present proteins it is most likely to bind.
  \choice It can be approached with multi-class classification.
  \choice It can be approached with multi-task classification.
  \choice For a given molecule, it determines every protein it would bind when several are present at once.
\end{itemize}
\closequestion
\begin{answer}{A, B}
  \yes{A} Deconvolution asks: of the candidate proteins, which is the \emph{most likely} target of this molecule.
  \yes{B} Picking one target out of several is naturally a \emph{multi-class} (one-of-$K$) problem.
  \no{C} \emph{Multi-task} classification predicts activity against \emph{many} targets independently --- that is target \emph{prediction}, not deconvolution.
  \no{D} Identifying \emph{all} binding proteins is again target prediction (many yes/no labels), not the ``single most likely'' framing of deconvolution.
\end{answer}

\questiontitle{Drug synergy and antagonism}
Which of the following statements about drug synergy and antagonism are correct?
\begin{itemize}[leftmargin=1.8em]
  \choice Synergy means the combined effect of two or more drugs exceeds the sum of their individual effects.
  \choice The Loewe additivity model is a standard tool for assessing synergy.
  \choice Antagonism means the combined effect of two or more drugs exceeds the sum of their individual effects.
  \choice Synergy/antagonism analysis applies only to chemically similar drugs.
\end{itemize}
\closequestion
\begin{answer}{A, B}
  \yes{A} Synergy = the combination does \emph{more} than the sum of the parts (super-additive).
  \yes{B} Loewe additivity is the standard reference model for ``no interaction'', against which synergy is measured.
  \no{C} This is the inverted definition: \emph{antagonism} is when the combined effect is \emph{less} than the sum. The sentence describes synergy but labels it antagonism.
  \no{D} Combinations are studied for \emph{any} drug pair (often deliberately different mechanisms) --- chemical similarity is irrelevant.
\end{answer}

\questiontitle{Random Forests vs.\ Graph Neural Networks}
Comparing Random Forests (RF) and Graph Neural Networks (GNNs) for molecular
modeling, which statements are correct?
\begin{itemize}[leftmargin=1.8em]
  \choice GNNs are inherently superior for QSAR because they capture the graph structure of molecules, giving more accurate predictions in every case.
  \choice RF models overfit less than GNNs, which makes them more reliable on datasets with many features but few samples.
  \choice RF models are usually more interpretable than GNNs, so the contribution of individual molecular features is easier to read off.
  \choice GNNs are less sensitive to input-data quality and variability than RFs, making them more consistent across different datasets.
\end{itemize}
\closequestion
\begin{answer}{B, C}
  \no{A} ``Better in \emph{all} cases'' is the giveaway --- on limited/noisy data RF often matches or beats GNNs (no free lunch).
  \yes{B} On small, high-dimensional QSAR datasets RF (bagging) is robust and overfits less than a data-hungry GNN.
  \yes{C} RF gives feature-importance scores over interpretable descriptors; GNN representations are opaque.
  \no{D} The opposite is true: GNNs are typically \emph{more} sensitive to data quantity/quality, not less.
\end{answer}


% ==========================================================
%  SECTION 3 — Graph Neural Networks & Advanced Methods
% ==========================================================
\section{Graph Neural Networks \& Advanced Methods}

\questiontitle{Molecular graphs and graph neural networks}
Which statements about molecular graphs and graph neural networks are correct?
\begin{itemize}[leftmargin=1.8em]
  \choice A molecular graph can be stored as two matrices: one for the node features and one for the edges between nodes.
  \choice Node features may include the atomic number, hybridization state, and formal charge.
  \choice A GNN uses only the node features to form its final prediction.
  \choice Bond-type information (e.g.\ single vs.\ double bond) cannot be represented in a GNN.
\end{itemize}
\closequestion
\begin{answer}{A, B}
  \yes{A} Standard encoding: a node-feature matrix plus an adjacency/edge matrix.
  \yes{B} Typical atom features are exactly atomic number, hybridisation, formal charge (one-hot/numeric).
  \no{C} Message passing also uses \emph{edge} information and the graph topology, not node features alone.
  \no{D} Bond type \emph{can} be encoded as edge features --- single/double/aromatic bonds are routinely included.
\end{answer}

\questiontitle{Limitations of GNNs in molecular modeling}
Which of the following are genuine limitations of GNNs in the context of
molecular modeling?
\begin{itemize}[leftmargin=1.8em]
  \choice GNNs cannot predict any molecular property with accuracy.
  \choice GNNs may struggle to capture long-range interactions between atoms.
  \choice GNNs can have trouble with very large, complex molecular graphs.
  \choice GNNs need considerable computational resources to train.
\end{itemize}
\closequestion
\begin{answer}{B, C, D}
  \no{A} Far too strong --- GNNs predict many molecular properties very well; the statement is an obvious overstatement.
  \yes{B} A fixed number of message-passing hops limits how far information travels (over-smoothing / under-reaching).
  \yes{C} Large, complex graphs strain both expressivity and memory.
  \yes{D} Training GNNs is compute- and memory-intensive relative to e.g.\ RF.
\end{answer}

\questiontitle{Multi-modal contrastive learning}
Which statements about multi-modal contrastive learning are correct?
\begin{itemize}[leftmargin=1.8em]
  \choice Multi-modal contrastive learning embeds several modalities (e.g.\ molecules and proteins) into one shared latent space.
  \choice Multi-modal contrastive learning is a self-supervised method.
  \choice Multi-modal contrastive learning concentrates the (latent) representations of several modalities (e.g.\ molecules and proteins) to make predictions.
  \choice Multi-modal contrastive learning is an unsupervised method.
\end{itemize}
\closequestion
\begin{answer}{A, B}
  \yes{A} The defining idea: map both modalities into one \emph{shared} embedding space (à la CLIP).
  \yes{B} It is \emph{self-supervised} --- supervision comes from matched/mismatched pairs, not human labels.
  \no{C} It \emph{aligns} matched pairs (pull together / push apart); ``concentrates representations to make predictions'' mischaracterises the objective.
  \no{D} The subtle distinction: it is self-supervised, \emph{not} unsupervised --- the pairing provides the training signal.
\end{answer}

\questiontitle{Multiple instance learning (MIL)}
Which of the following statements about multiple instance learning (MIL) hold?
\begin{itemize}[leftmargin=1.8em]
  \choice MIL trains on sets (bags) of instances in which each instance carries its own label.
  \choice MIL needs the bag size to be fixed.
  \choice MIL trains on sets (bags) of instances that share a single label.
  \choice A MIL model should be invariant to the ordering of the instances.
\end{itemize}
\closequestion
\begin{answer}{C, D}
  \no{A} Exactly what MIL is \emph{not}: instances are \emph{not} individually labelled (that would be ordinary supervised learning).
  \no{B} Bags can have varying sizes; the aggregation (e.g.\ attention pooling) handles any number of instances.
  \yes{C} The hallmark of MIL: a label is attached to the \emph{bag}, not to each instance.
  \yes{D} A bag is a set, so the model output must be \emph{permutation-invariant}.
\end{answer}


% ==========================================================
%  SECTION 4 — Generative Models & Chemical Synthesis
% ==========================================================
\section{Generative Models \& Chemical Synthesis}

\questiontitle{Generative models for small molecules}
Which statements about generative models for small molecules are correct?
\begin{itemize}[leftmargin=1.8em]
  \choice Because SMILES syntax is intricate with long-range dependencies, SMILES generators yield fewer than 10\% valid strings.
  \choice Generative models are used to explore the already-synthesised, known chemical space.
  \choice Generative adversarial networks cannot generate small molecules.
  \choice Autoregressive models can generate SMILES representations of small molecules.
\end{itemize}
\closequestion
\begin{answer}{D}
  \no{A} Empirically false: an LSTM reaches $\sim$98\% valid SMILES, and SELFIES gives 100\% by construction --- nowhere near ``$<$10\%''.
  \no{B} Generative models explore \emph{novel} chemical space, not just the already-synthesised \emph{known} set.
  \no{C} GANs \emph{can} generate molecules (as can VAEs, normalizing flows, diffusion).
  \yes{D} Autoregressive (next-token) models generate SMILES one character at a time and work well.
\end{answer}

\questiontitle{Goal-directed molecule generation}
Which of the following correctly describe goal-directed molecule generation?
\begin{itemize}[leftmargin=1.8em]
  \choice Goal-directed generation is concerned only with reproducing existing structures that already have the desired properties.
  \choice Goal-directed generation aims to produce molecules whose properties resemble those of the training set.
  \choice Reinforcement learning can be used to drive goal-directed molecule generation.
  \choice Goal-directed generation aims to produce molecules with particular, pre-specified properties.
\end{itemize}
\closequestion
\begin{answer}{C, D}
  \no{A} ``Replicating existing structures'' describes \emph{distribution learning}, not goal-directed design.
  \no{B} ``Similar to the training set'' is again \emph{distribution learning} --- goal-directed deliberately pushes \emph{beyond} the training distribution.
  \yes{C} A scoring function + RL loop (generate, score, reward, update) is a standard goal-directed recipe (e.g.\ REINVENT).
  \yes{D} The aim is molecules with \emph{specific desired} properties, optimised toward an objective.
\end{answer}

\questiontitle{Fr\'echet ChemNet Distance (FCD)}
Which statements about the Fr\'echet ChemNet Distance (FCD) and how it is computed
are correct?
\begin{itemize}[leftmargin=1.8em]
  \choice FCD measures the distance between the individual atoms of molecules.
  \choice FCD compares the means and covariances of the molecules' latent representations.
  \choice FCD is computed straight from the raw molecular structures with no transformation.
  \choice FCD is a metric comparing the distribution of generated molecules with that of real ones.
\end{itemize}
\closequestion
\begin{answer}{B, D}
  \no{A} It works on whole-molecule embeddings, not distances between individual atoms.
  \yes{B} FCD fits Gaussians to ChemNet's last-hidden-layer embeddings and compares their \emph{means and covariances}.
  \no{C} It is computed in a \emph{learned latent space} (ChemNet), not on raw structures --- there is a transformation step.
  \yes{D} It is a \emph{distributional} metric: how close is the generated set to the real set overall.
\end{answer}

\questiontitle{AI for chemical reactions}
Which of the following statements about machine learning for chemical reactions are
correct?
\begin{itemize}[leftmargin=1.8em]
  \choice Some approaches use so-called reaction templates, i.e.\ transformation rules for molecules.
  \choice Deep-learning reaction-outcome predictors can be trained on databases of known reactions.
  \choice Transformer architectures have been applied to predicting reaction outcomes.
  \choice Predicting reactions with machine learning requires quantum-mechanical data because of bond changes and electron movement.
\end{itemize}
\closequestion
\begin{answer}{A, B, C}
  \yes{A} Template-based (neural-symbolic) methods select a reaction template = a transformation rule.
  \yes{B} Models are trained on reaction databases like USPTO (mined from patents).
  \yes{C} The ``Molecular Transformer'' treats reaction prediction as seq-to-seq translation.
  \no{D} No QM data needed --- predictions come from \emph{reaction databases}, not quantum-mechanical calculations.
\end{answer}

\questiontitle{Synthesis planning}
Which statements about synthesis planning are correct?
\begin{itemize}[leftmargin=1.8em]
  \choice Synthesis planning designs a sequence of chemical reactions that yields a target compound.
  \choice Retrosynthetic analysis is central to it, decomposing complex molecules into simpler starting materials.
  \choice Modern deep-learning methods have made synthesis planning trivial.
  \choice Synthesis planning amounts to nothing more than choosing the reactants.
\end{itemize}
\closequestion
\begin{answer}{A, B}
  \yes{A} It designs a whole \emph{sequence} of reactions from purchasable materials to the target.
  \yes{B} Retrosynthesis recursively breaks the target into simpler precursors --- the core technique.
  \no{C} Far from trivial: $\sim$50\% solve rates, and whether routes actually run under real conditions remains uncertain.
  \no{D} It is much more than ``selecting reactants'' --- it is multi-step route design with scoring and search.
\end{answer}

\questiontitle{Retrosynthesis as a reinforcement-learning game}
Multi-step retrosynthesis can be cast as a game solvable by reinforcement learning.
Which of the following statements about this ``retrosynthesis game'' are correct?
\begin{itemize}[leftmargin=1.8em]
  \choice In the ``game'', an action consists of choosing reactants for a reaction.
  \choice In the ``game'', the state is the set of available building blocks.
  \choice In the ``game'', the environment carries out the reaction and returns the new molecules.
  \choice In the ``game'', the reward is granted once the product has been synthesised from the building blocks.
\end{itemize}
\closequestion
\begin{answer}{A, C, D}
  \yes{A} An \emph{action} is choosing a disconnection / set of reactants (a reverse reaction).
  \no{B} The classic trap: the \emph{state} is the current set of molecules \emph{still to be made}, \emph{not} the set of available building blocks (those are the goal you are trying to reach).
  \yes{C} The \emph{environment} applies the chosen reaction and returns the new set of molecules.
  \yes{D} The \emph{reward} comes when every remaining molecule is an available building block (the game is won).
\end{answer}


% ==========================================================
%  SECTION 5 — Medical Data & Imaging
% ==========================================================
\section{Medical Data \& Imaging}

\questiontitle{Tabular medical data: types and missing values}
Suppose we are given a table of patient medical records with three columns:
``smoking status'' (one of three string values: non-smoker, former smoker,
current smoker), ``cholesterol level'' (floating-point measurements in mg/dL), and
``diabetes risk score'' (integer values in $[0,10]$ indicating the likelihood of
developing diabetes). Which of the following statements are correct?
\begin{itemize}[leftmargin=1.8em]
  \choice Missing ``cholesterol level'' entries can be filled with the column's median to cope with skew.
  \choice The ``smoking status'' column is continuous, since it reflects patients' various smoking habits.
  \choice ``cholesterol level'' is ordinal, because the measurements can be placed in increasing order.
  \choice Missing ``smoking status'' entries can be filled with the column's mean to cope with skew.
\end{itemize}
\closequestion
\begin{answer}{A}
  \yes{A} Cholesterol is numeric; the \emph{median} is the robust choice for imputing a skewed continuous column.
  \no{B} Smoking status is \emph{categorical} (three discrete labels), not continuous.
  \no{C} Being orderable is not enough: cholesterol is a true \emph{continuous} (interval/ratio) measurement with meaningful differences, not merely ordinal.
  \no{D} You cannot take a \emph{mean} of a categorical column --- a missing smoking status would be imputed with the \emph{mode}, not the mean.
\end{answer}

\questiontitle{Batch effects in medical data}
Which statements about batch effects in medical data are correct?
\begin{itemize}[leftmargin=1.8em]
  \choice Overlooking batch effects can yield models that separate batches instead of the true biological conditions.
  \choice Batch effects vanish simply by collecting more samples.
  \choice Batch effects introduce artificial differences that stem from varying experimental conditions.
  \choice Batch effects are generally helpful, supplying needed variability to the data.
\end{itemize}
\closequestion
\begin{answer}{A, C}
  \yes{A} If batches correlate with the label, the model takes the shortcut and separates \emph{batches} instead of biology.
  \no{B} More samples from the \emph{same} biased batches does not remove the systematic offset --- you need design/correction (randomisation, ComBat, etc.).
  \yes{C} They are exactly artificial, non-biological differences from instruments/reagents/handling.
  \no{D} They are a confounder, not a benefit --- this is the opposite of correct.
\end{answer}

\questiontitle{Domain shifts in medical and healthcare data}
Medical and healthcare data are frequently subject to domain shifts. Using the
lecture notation ($y$ = disease status, e.g.\ ``healthy''/``diseased'';
$x$ = patient features, e.g.\ age, blood parameters, heart rate), which of the
following statements about domain shifts are correct?
\begin{itemize}[leftmargin=1.8em]
  \choice A model with high predictive quality on a randomly drawn validation set will, by that fact alone, adapt to domain shifts and keep its quality over time.
  \choice The patient-feature distribution $p(x)$ is never touched by domain shifts.
  \choice During waves of an infectious disease (such as COVID-19), the prevalence---and thus the probability $p(y)$ of observing the phenotype---changes.
  \choice If the way the status $y$ is diagnosed changes, e.g.\ an antigen test is replaced by a PCR test, this is a concept shift $p(y\mid x)$.
\end{itemize}
\closequestion
\begin{answer}{C, D}
  \no{A} A good \emph{random} validation score says nothing about future shifts; models do \emph{not} auto-adapt over time.
  \no{B} ``Never'' is wrong --- $p(x)$ \emph{does} shift (covariate shift), e.g.\ a different patient population.
  \yes{C} A disease wave changes the prevalence $p(y)$ --- a prior/label shift.
  \yes{D} Redefining how the label is measured changes the $x\to y$ relationship: a concept shift $p(y\mid x)$.
\end{answer}

\questiontitle{Medical imaging modalities}
Which statements about medical imaging modalities are correct?
\begin{itemize}[leftmargin=1.8em]
  \choice Computed Tomography (CT) builds its body images from ultrasound.
  \choice X-rays distinguish tissues of differing density, such as bone versus soft tissue.
  \choice Magnetic Resonance Imaging (MRI) yields detailed images of internal structures using strong magnetic fields and radio waves.
  \choice Optical coherence tomography forms images of organs and tissues using multiple X-rays.
\end{itemize}
\closequestion
\begin{answer}{B, C}
  \no{A} CT is built from many \emph{X-ray} projections, not ultrasound.
  \yes{B} X-ray contrast comes from differential absorption by tissue density (bone vs.\ soft tissue).
  \yes{C} MRI uses strong magnetic fields + radio-frequency pulses (no ionising radiation).
  \no{D} OCT uses \emph{light} (near-infrared interferometry), not X-rays.
\end{answer}

\questiontitle{Most common task in medical imaging}
Compared with typical computer-vision tasks, which task is especially prominent in
medical imaging?
\begin{itemize}[leftmargin=1.8em]
  \choice Image generation.
  \choice Classification.
  \choice Segmentation.
  \choice Object detection.
\end{itemize}
\closequestion
\begin{answer}{B, C}
  \no{A} Image generation is a niche/auxiliary use, not the dominant clinical task.
  \yes{B} Classification (e.g.\ disease present/absent, grading) is also very common.
  \yes{C} \emph{Segmentation} (delineating organs/lesions pixel-by-pixel) is the signature, disproportionately common medical-imaging task --- hence the U-Net.
  \no{D} Object detection is comparatively rare in medical imaging versus segmentation and classification.
\end{answer}

\questiontitle{The U-Net architecture}
Which of the following correctly describe the U-Net architecture?
\begin{itemize}[leftmargin=1.8em]
  \choice U-Net was the first vision model whose encoder relied on self-attention rather than convolutions.
  \choice The U-Net architecture has a symmetric encoder and decoder.
  \choice U-Net was created specifically for the object-detection task.
  \choice To help decoding, encoder feature maps are concatenated with those of the matching up-sampling level of the decoder.
\end{itemize}
\closequestion
\begin{answer}{B, D}
  \no{A} U-Net is fully \emph{convolutional}; it predates and does not use self-attention encoders.
  \yes{B} The ``U'' is a symmetric contracting encoder + expanding decoder.
  \no{C} It was designed for \emph{segmentation} (biomedical), not object detection.
  \yes{D} Skip connections concatenate encoder feature maps into the matching decoder level --- the key idea that preserves spatial detail.
\end{answer}

\questiontitle{The ``Clever Hans'' effect}
In the context of AI for medical imaging, what does the ``Clever Hans'' effect
refer to?
\begin{itemize}[leftmargin=1.8em]
  \choice A bias that arises when up-scaling the resolution of low-quality medical images.
  \choice A case where a model looks like it performs well but is actually leaning on irrelevant features or artifacts in the data.
  \choice A case where an AI predicts outcomes incorrectly because its training data were biased.
  \choice The use of data augmentation to make medical-imaging models more robust.
\end{itemize}
\closequestion
\begin{answer}{B}
  \no{A} Unrelated to super-resolution bias.
  \yes{B} Exactly shortcut learning: the model \emph{looks} accurate but is exploiting an irrelevant artifact (e.g.\ a hospital tag/ruler), not the pathology.
  \no{C} Tempting, but too generic: ``biased training data'' describes \emph{any} bias. Clever Hans specifically means relying on \emph{spurious shortcut features} while appearing to perform well.
  \no{D} Data augmentation is a mitigation technique, not the Clever Hans effect.
\end{answer}

\questiontitle{Explainable AI and interpretability in medical imaging}
Which statements about explainable AI and interpretability in medical imaging are
correct?
\begin{itemize}[leftmargin=1.8em]
  \choice Gradient integration propagates a relevance value backwards through the layers to the input, estimating each input feature's relevance.
  \choice Explainable-AI methods such as Grad-CAM help visualise which parts of a medical image matter most for a prediction.
  \choice Occlusion analysis deletes patches from the training images to make the model more interpretable.
  \choice Counterfactual explanations generate alternative images showing what would have had to change for a different diagnosis.
\end{itemize}
\closequestion
\begin{answer}{B, D}
  \no{A} The layer-by-layer back-propagation of a \emph{relevance} value is Layer-wise Relevance Propagation (LRP); ``gradient integration'' is the wrong name for it.
  \yes{B} Grad-CAM produces a heatmap of the image regions most responsible for the prediction.
  \no{C} Occlusion perturbs the \emph{input} image at inference (mask a patch, see the score drop) --- it does \emph{not} remove training data, and it explains rather than ``improves interpretability''.
  \yes{D} Counterfactuals show the minimal image change that would flip the diagnosis --- a what-if explanation.
\end{answer}

\questiontitle{AlphaFold 2}
Which of the following statements about AlphaFold 2 are correct?
\begin{itemize}[leftmargin=1.8em]
  \choice AlphaFold 2 relies on quantum-mechanical calculations to sharpen its predictions.
  \choice AlphaFold 2 draws on multiple sequence alignments and homologous protein structures to improve its accuracy.
  \choice AlphaFold 2 is a deep-learning model that predicts a protein's three-dimensional structure from its amino-acid sequence.
  \choice AlphaFold 2 is built on a convolutional neural network architecture.
\end{itemize}
\closequestion
\begin{answer}{B, C}
  \no{A} No quantum-mechanical calculations are used --- it is a learned model.
  \yes{B} It relies heavily on multiple sequence alignments (co-evolution signal) and template structures.
  \yes{C} Its job is exactly sequence $\to$ 3D structure prediction.
  \no{D} The architecture is \emph{attention}-based (the Evoformer + structure module), not a CNN.
\end{answer}


% ==========================================================
%  SECTION 6 — Biological Sequences
% ==========================================================
\section{Biological Sequences}

\questiontitle{Composition of biological sequences}
Which statements about the composition of biological sequences are correct?
\begin{itemize}[leftmargin=1.8em]
  \choice Protein sequences are chains of amino acids joined by peptide bonds.
  \choice DNA sequences are built from nucleic acids.
  \choice Biological sequences obey no particular patterns or rules in their make-up.
  \choice RNA sequences are built from lipids.
\end{itemize}
\closequestion
\begin{answer}{A, B}
  \yes{A} Proteins are amino-acid chains joined by peptide bonds.
  \yes{B} DNA is a chain of nucleotides (nucleic acid).
  \no{C} Sequences are highly structured (motifs, codons, grammar) --- ``no rules'' is false and is what makes them learnable.
  \no{D} RNA is a \emph{nucleic acid}, not a lipid.
\end{answer}

\questiontitle{Deep learning for biological sequences}
Which of the following statements about deep learning methods for biological
sequences are correct?
\begin{itemize}[leftmargin=1.8em]
  \choice Feature extraction from biological sequences can reveal motifs, i.e.\ recurring patterns or elements within them.
  \choice One-hot encoding is a common way to prepare biological sequences as input for deep models.
  \choice Deep-learning techniques apply to tasks such as classification, segmentation, and generative modelling of biological sequences.
  \choice 1D convolutions, Recurrent Neural Networks (RNNs), Long Short-Term Memory networks (LSTMs), and Transformers are all used to analyse biological sequences.
\end{itemize}
\closequestion
\begin{answer}{A, B, C, D}
  \yes{A} Learned filters detect recurring \emph{motifs} (e.g.\ binding sites).
  \yes{B} One-hot encoding of the alphabet (A/C/G/T or 20 amino acids) is the standard input format.
  \yes{C} Classification, segmentation (per-position labels) and generation are all applied to sequences.
  \yes{D} 1D-CNNs, RNNs, LSTMs and Transformers are all standard sequence architectures here --- every option is correct.
\end{answer}


% ==========================================================
%  2026 EXAM (single-best-answer) — Sections 7–10
% ==========================================================
\begin{notebox}[2026 exam --- single-best-answer]
The remaining sections are from the \textbf{2026} exam (25 June 2026). Each question
has \textbf{exactly one} correct option among the three listed; the answer box marks
it \textcolor{ok}{\ensuremath{\checkmark}} and explains why the other two are wrong.
\end{notebox}


% ==========================================================
%  SECTION 7 — Molecules, Representations & Cheminformatics (2026)
% ==========================================================
\section{Molecules, Representations \& Cheminformatics (2026)}

\questiontitle{Molecular graphs vs.\ fingerprints}
Which of the following statements accurately describes the relationship between
molecular graphs and molecular fingerprints?
\begin{itemize}[leftmargin=1.8em]
  \choice A fingerprint can be computed from a molecular graph, yet the exact graph usually cannot be recovered from the fingerprint.
  \choice Fingerprints and molecular graphs must hold the same structural information.
  \choice A molecular graph can always be reconstructed exactly from its fingerprint.
\end{itemize}
\closequestion
\begin{answer}{A}
  \yes{A} A fingerprint is \emph{computed from} the graph (substructure enumeration/hashing), but the mapping is lossy and many-to-one, so the exact graph cannot be recovered.
  \no{B} A fingerprint is a compressed bit/count summary that discards information, so it does \emph{not} carry the same structural information as the graph.
  \no{C} Exact reconstruction is impossible in general --- different molecules can even hash to the same bits.
\end{answer}

\questiontitle{Converting between SMILES and fingerprints}
Using standard cheminformatics tools, in which direction is conversion between a
SMILES string and a molecular fingerprint actually possible?
\begin{itemize}[leftmargin=1.8em]
  \choice Cheminformatics tools can convert SMILES and fingerprints into one another in both directions.
  \choice Cheminformatics tools can turn a SMILES into a fingerprint, but not the other way round.
  \choice Cheminformatics tools can turn a fingerprint into a SMILES, but not the other way round.
\end{itemize}
\closequestion
\begin{answer}{B}
  \no{A} The round-trip does not exist: you cannot go fingerprint $\to$ SMILES.
  \yes{B} SMILES $\to$ molecular graph $\to$ fingerprint is a standard cheminformatics pipeline.
  \no{C} Backwards --- fingerprints are lossy, so fingerprint $\to$ SMILES is not possible.
\end{answer}

\questiontitle{Tanimoto coefficient}
On what quantity is the Tanimoto coefficient based when it scores how similar two
molecular fingerprints are?
\begin{itemize}[leftmargin=1.8em]
  \choice It takes the count of shared features over the total number of features present in either molecule.
  \choice It superimposes three-dimensional pharmacophore points and averages their spatial distances.
  \choice It compares the core scaffolds, scoring one whenever the scaffolds coincide.
\end{itemize}
\closequestion
\begin{answer}{A}
  \yes{A} Tanimoto $= |A\cap B| / |A\cup B|$: shared on-bits divided by the union of on-bits.
  \no{B} Averaging 3D pharmacophore-point distances is a different (3D alignment) similarity, not Tanimoto.
  \no{C} Tanimoto does not test scaffold identity; it counts shared fingerprint features.
\end{answer}

\questiontitle{Aspirin SMILES}
Which of the following SMILES strings represents aspirin (acetylsalicylic acid)?
\begin{itemize}[leftmargin=1.8em]
  \choice \texttt{O=C(C)Oc1ccccc1C(=O)O}
  \choice \texttt{cc(=O)Oc1ccccc1c(=O)O}
  \choice \texttt{O=C[C]Oc1ccccc1C[=O]O}
\end{itemize}
\closequestion
\begin{answer}{A}
  \yes{A} \texttt{O=C(C)Oc1ccccc1C(=O)O}: an acetyl ester (\texttt{O=C(C)O}) on an aromatic ring (\texttt{c1ccccc1}) bearing a carboxylic acid (\texttt{C(=O)O}) --- acetylsalicylic acid.
  \no{B} Lowercase \texttt{c} marks \emph{aromatic} atoms, so writing the carbonyl carbons as aromatic (\texttt{cc(=O)}) is chemically and syntactically wrong.
  \no{C} \texttt{[C]} and \texttt{[=O]} are malformed SMILES atom tokens --- the string does not parse.
\end{answer}

\questiontitle{Bioactivity datasets}
Which of the following best captures what bioactivity datasets usually look like in practice?
\begin{itemize}[leftmargin=1.8em]
  \choice They are usually large, balanced, and almost free of measurement noise.
  \choice They are frequently small, imbalanced, and noisy, with far more inactive than active compounds.
  \choice They are cheap to produce and therefore cover most compound--assay combinations.
\end{itemize}
\closequestion
\begin{answer}{B}
  \no{A} The opposite of reality --- bioactivity data is rarely large, balanced, or clean.
  \yes{B} It is typically limited, imbalanced (far more inactives than actives), and noisy.
  \no{C} Assays are expensive, so coverage of compound$\times$assay combinations is very sparse.
\end{answer}

\questiontitle{The Design--Make--Test--Analyse cycle}
In drug discovery, which option correctly captures how the Design--Make--Test--Analyse
(DMTA) cycle proceeds?
\begin{itemize}[leftmargin=1.8em]
  \choice Compounds are tested before being designed, and only the active ones are afterwards synthesized.
  \choice Compounds are designed and made just once, after which testing removes any need for further design.
  \choice Compounds are designed, made, and tested experimentally, and the analysis steers the next design round.
\end{itemize}
\closequestion
\begin{answer}{C}
  \no{A} Designing comes before testing --- you cannot test a compound you have not yet designed and made.
  \no{B} It is an iterative \emph{cycle}; one round of testing does not remove the need for further design.
  \yes{C} Design $\to$ Make $\to$ Test $\to$ Analyse, and the analysis feeds the next design iteration.
\end{answer}


% ==========================================================
%  SECTION 8 — Virtual Screening, QSAR & Targets (2026)
% ==========================================================
\section{Virtual Screening, QSAR \& Targets (2026)}

\questiontitle{Structure-based vs.\ ligand-based drug design}
What is the essential difference between structure-based drug design (SBDD) and
ligand-based drug design (LBDD)?
\begin{itemize}[leftmargin=1.8em]
  \choice SBDD needs known active ligands, whereas LBDD needs the target's 3D structure.
  \choice Both SBDD and LBDD need a high-resolution 3D structure of the target.
  \choice SBDD relies on the target's 3D structure, whereas LBDD can work from known ligands without that structure.
\end{itemize}
\closequestion
\begin{answer}{C}
  \no{A} Swapped: SBDD needs the target structure, LBDD needs known ligands.
  \no{B} LBDD specifically does \emph{not} require the target's 3D structure.
  \yes{C} SBDD uses the target's 3D structure; LBDD relies on known ligands without it.
\end{answer}

\questiontitle{Similarity search vs.\ model-based search}
In what way do similarity search and model-based search differ from each other?
\begin{itemize}[leftmargin=1.8em]
  \choice Similarity search is ligand-based, whereas model-based search must have the 3D structure of the target's binding site.
  \choice Similarity search retrieves molecules resembling a query, whereas model-based search predicts activity with a model trained on labelled data.
  \choice Similarity search learns a decision function from labelled actives and inactives, whereas model-based search ranks compounds by closeness to a query.
\end{itemize}
\closequestion
\begin{answer}{B}
  \no{A} Model-based search does \emph{not} necessarily need the target's 3D structure --- it can be a ligand-based supervised model.
  \yes{B} Similarity search retrieves neighbours of a query molecule; model-based search predicts activity with a model trained on labels.
  \no{C} The two definitions are reversed.
\end{answer}

\questiontitle{Early-recognition metrics (odd one out)}
Of the metrics below, which one is \emph{not} built to reward finding actives early,
i.e.\ near the top of a ranked virtual-screening list?
\begin{itemize}[leftmargin=1.8em]
  \choice BEDROC.
  \choice AUC-ROC.
  \choice Enrichment factor.
\end{itemize}
\closequestion
\begin{answer}{B}
  \no{A} BEDROC is built specifically to reward early recognition.
  \yes{B} AUC-ROC weights all ranking positions equally, so it is \emph{not} an early-recognition metric --- the odd one out.
  \no{C} The enrichment factor measures actives in the top $x\%$ --- an early-recognition metric.
\end{answer}

\questiontitle{Evaluating a virtual-screening model}
A virtual-screening model orders a million compounds, yet your budget only permits
assaying the top 1\%. Which model-selection strategy best fits this situation?
\begin{itemize}[leftmargin=1.8em]
  \choice Pick the model with the best accuracy at a 0.5 probability threshold, since ranking above that threshold no longer matters once compounds are classified.
  \choice Pick the model with the highest overall AUC-ROC, since weighting all ranking positions equally guarantees the best enrichment in the top 1\%.
  \choice Favour an early-recognition metric such as the enrichment factor near the 1\% cutoff or BEDROC, since the top of the ranking decides experimental usefulness.
\end{itemize}
\closequestion
\begin{answer}{C}
  \no{A} A 0.5-threshold accuracy ignores how well the very top of the ranking is ordered --- which is all that matters here.
  \no{B} Global AUC-ROC does \emph{not} guarantee good enrichment inside the top 1\%.
  \yes{C} Use an early-recognition metric --- enrichment factor near the 1\% cutoff, or BEDROC.
\end{answer}

\questiontitle{Defensible generalization estimates in QSAR}
In a QSAR pipeline, which procedure yields the most trustworthy estimate of how well
the model generalises?
\begin{itemize}[leftmargin=1.8em]
  \choice Fix the split first, carry out feature and hyperparameter selection on the training data, and reserve the held-out test set for the final evaluation only.
  \choice Split first, then compare models and hyperparameters on the test set and report the best test-set score.
  \choice Select features on the whole dataset, then split the reduced data and tune hyperparameters on the training part only.
\end{itemize}
\closequestion
\begin{answer}{A}
  \yes{A} Split first, do all feature/hyperparameter selection on the training data, and touch the test set only once for the final estimate.
  \no{B} Tuning on the test set leaks it into model selection --- the reported number is optimistic.
  \no{C} Selecting features on the full dataset before splitting also leaks test information.
\end{answer}

\questiontitle{Compound-series bias}
What does compound-series bias refer to, and why does it undermine QSAR model evaluation?
\begin{itemize}[leftmargin=1.8em]
  \choice Because many compounds are closely related, a random split can put near-analogues in both training and test, underestimating the generalization error.
  \choice Test-compound information leaks into feature or hyperparameter selection, biasing the reported score even when compounds are structurally unrelated.
  \choice The active class is far smaller than the inactive one, so the majority class dominates performance even when the chemical series are diverse.
\end{itemize}
\closequestion
\begin{answer}{A}
  \yes{A} Close analogues land in both train and test under a random split, so the test set is partly memorised and generalization error is underestimated.
  \no{B} That describes generic information leakage in feature/hyperparameter selection, not compound-series bias.
  \no{C} That describes class imbalance, a different problem.
\end{answer}

\questiontitle{Drug synergy and model symmetry}
Which option combines a correct definition of drug synergy with a sensible property
that a model predicting it should have?
\begin{itemize}[leftmargin=1.8em]
  \choice The combined effect is larger than the additive one, and the prediction should stay the same when the two compounds are swapped.
  \choice The combined effect equals the additive expectation, and the prediction may differ when the pair is supplied in reverse order.
  \choice The combined effect is below the additive expectation, and the prediction should be unaffected by the order of the compounds.
\end{itemize}
\closequestion
\begin{answer}{A}
  \yes{A} Synergy is super-additive (combined $>$ additive), and a sensible model is symmetric --- swapping the two compounds must not change the prediction.
  \no{B} ``Equals the additive expectation'' is \emph{no} interaction, and order-dependence is an undesirable property.
  \no{C} ``Smaller than additive'' is \emph{antagonism}, not synergy (the symmetry property is right, but the definition is wrong).
\end{answer}

\questiontitle{Proteochemometric modelling}
In what respect does proteochemometric modelling go beyond classical single-target QSAR?
\begin{itemize}[leftmargin=1.8em]
  \choice It uses representations of both the molecule and the protein target together.
  \choice It uses molecular representations with a softmax over the candidate protein targets.
  \choice It uses molecular representations with one separate prediction head per protein target.
\end{itemize}
\closequestion
\begin{answer}{A}
  \yes{A} Proteochemometrics jointly encodes the \emph{molecule and the protein target}, so one model spans many targets.
  \no{B} A softmax over targets from the molecule alone has no target representation --- not the proteochemometric idea.
  \no{C} A separate head per target (multi-task) still uses only molecular inputs --- not a joint molecule+protein representation.
\end{answer}

\questiontitle{Target deconvolution / target fishing}
How is the task known as target deconvolution (or target fishing) defined?
\begin{itemize}[leftmargin=1.8em]
  \choice Given one protein and a compound library, rank every molecule by predicted affinity to that protein.
  \choice Given one molecule and several candidate proteins, predict which protein it is most likely to bind.
  \choice Given one molecule, predict binding to each candidate protein independently, so any number of targets may be assigned at once.
\end{itemize}
\closequestion
\begin{answer}{B}
  \no{A} Ranking a whole library against \emph{one} protein is virtual screening --- the inverse problem.
  \yes{B} Target fishing: given a molecule and several candidate proteins, identify the \emph{most likely} target.
  \no{C} Predicting binding to \emph{every} target independently is multi-task target prediction, not deconvolution's ``single most likely target'' framing.
\end{answer}


% ==========================================================
%  SECTION 9 — Graph Neural Networks (2026)
% ==========================================================
\section{Graph Neural Networks (2026)}

\questiontitle{Reordering atoms in a molecular graph}
Suppose the atoms of a molecular graph are given new indices while the bonds between
them are left unchanged. What behaviour do we expect from a graph neural network?
\begin{itemize}[leftmargin=1.8em]
  \choice The node representations are permuted accordingly, but the graph-level prediction may change because the adjacency matrix was reordered.
  \choice Both the node representations and the graph-level prediction stay in exactly the same tensor positions.
  \choice The node representations are permuted accordingly, while the graph-level prediction stays unchanged.
\end{itemize}
\closequestion
\begin{answer}{C}
  \no{A} A pure relabelling must \emph{not} change the graph-level prediction --- it is the same graph, just reindexed.
  \no{B} Node representations should follow the permutation, not stay in fixed positions.
  \yes{C} Correct symmetry: node features are \emph{equivariant} (they permute with the atoms) while the graph-level readout is \emph{invariant}.
\end{answer}

\questiontitle{Steps of an MPNN layer}
What is the usual order of operations carried out inside a single Message Passing
Neural Network (MPNN) layer?
\begin{itemize}[leftmargin=1.8em]
  \choice Generate messages, aggregate them, then update.
  \choice Generate messages, apply graph convolution, then aggregate.
  \choice Apply graph convolution, generate messages, then update.
\end{itemize}
\closequestion
\begin{answer}{A}
  \yes{A} One MPNN layer: generate messages, aggregate them at each node, then update the node state.
  \no{B} ``Graph convolution'' is not a separate sub-step here, and aggregation is not the final step.
  \no{C} Wrong order and wrong steps --- the update must come last.
\end{answer}

\questiontitle{Oversmoothing in GNNs}
Once a large number of message-passing layers are stacked, the node embeddings across
a molecular graph end up nearly identical. Why does this cause trouble?
\begin{itemize}[leftmargin=1.8em]
  \choice Long-range information is squeezed through narrow graph bottlenecks, so distant signals cannot travel effectively.
  \choice The model may no longer tell local atom environments apart because the node embeddings grow too similar.
  \choice Important information never reaches distant nodes because too few message-passing steps were taken.
\end{itemize}
\closequestion
\begin{answer}{B}
  \no{A} Distant signals failing through bottlenecks is \emph{over-squashing}, a different pathology.
  \yes{B} Oversmoothing: after many layers the node embeddings converge, so the model can no longer tell local atom environments apart.
  \no{C} ``Too few steps'' (under-reaching) is the opposite of the ``many layers'' premise.
\end{answer}

\questiontitle{Contrastive learning}
By what mechanism does contrastive learning build a joint molecule--target embedding
space?
\begin{itemize}[leftmargin=1.8em]
  \choice It clusters molecules and targets separately and assumes equal-index clusters are binding pairs.
  \choice It pulls associated molecule-target pairs closer together and pushes mismatched pairs apart.
  \choice It concatenates molecule and target representations and fits a supervised activity regressor without comparing pairs in a shared space.
\end{itemize}
\closequestion
\begin{answer}{B}
  \no{A} Clustering each modality separately and matching by cluster index is not contrastive learning.
  \yes{B} Contrastive learning pulls associated molecule--target pairs together and pushes mismatched pairs apart in a shared space.
  \no{C} Concatenating and fitting a supervised regressor never builds a shared comparison space.
\end{answer}

\questiontitle{Multiple Instance Learning}
How does Multiple Instance Learning depart from ordinary supervised learning, in
which every feature vector carries a single label $(\boldsymbol{x}, y)$?
\begin{itemize}[leftmargin=1.8em]
  \choice \ldots we are given a bag of feature vectors that share a single label $(\{\boldsymbol{x}_1, \ldots, \boldsymbol{x}_N\}, y)$.
  \choice \ldots we are given a bag of feature vectors paired with a bag of labels $(\{\boldsymbol{x}_1, \ldots, \boldsymbol{x}_N\}, \{y_1, \ldots, y_M\})$.
  \choice \ldots we are given single feature vectors paired with a bag of labels $(\boldsymbol{x}, \{y_1, \ldots, y_M\})$.
\end{itemize}
\closequestion
\begin{answer}{A}
  \yes{A} MIL replaces $(\boldsymbol{x}, y)$ with a \emph{bag} of vectors sharing one label, $(\{\boldsymbol{x}_1,\ldots,\boldsymbol{x}_N\}, y)$.
  \no{B} A bag labelled with a \emph{bag} of values is not MIL (that would be a multi-output/multi-label setup).
  \no{C} A single vector with multiple labels is ordinary multi-label learning, not MIL.
\end{answer}


% ==========================================================
%  SECTION 10 — Generative Models & Synthesis Planning (2026)
% ==========================================================
\section{Generative Models \& Synthesis Planning (2026)}

\questiontitle{Reward hacking a QSAR proxy}
A generator is trained to maximise a QSAR-based activity score and learns to output
molecules the score rates as highly active, yet later wet-lab assays find them
inactive. What most plausibly went wrong?
\begin{itemize}[leftmargin=1.8em]
  \choice The activity score was supplied only after each molecule was finished, so the reward was delayed.
  \choice The generator gamed weaknesses in the proxy scoring model instead of raising genuine experimental activity.
  \choice The generator made too few structurally distinct molecules, lowering the diversity of candidates.
\end{itemize}
\closequestion
\begin{answer}{B}
  \no{A} Delayed (sparse) reward does not explain why \emph{predicted} activity is high while \emph{real} activity is low.
  \yes{B} Classic reward hacking: the generator found inputs that fool the proxy QSAR model instead of being genuinely active.
  \no{C} Low diversity would not systematically inflate predicted activity above true activity.
\end{answer}

\questiontitle{Combating analogue rediscovery (odd one out)}
A goal-directed generator keeps producing close analogues of molecules it has already
scored highly. Which of the following would \emph{not} help to counteract this
tendency?
\begin{itemize}[leftmargin=1.8em]
  \choice Blend samples from a fixed exploratory generator with those of the trainable generator while generating SMILES.
  \choice Penalise the reward of molecules resembling earlier high-scoring ones kept in a memory buffer.
  \choice Push exploitation harder by always taking the highest-scoring molecule as the next training example.
\end{itemize}
\closequestion
\begin{answer}{C}
  \no{A} Mixing in a fixed exploratory generator adds diversity --- a genuine remedy, so not the answer.
  \no{B} Penalising similarity to a memory of past hits also promotes diversity --- a genuine remedy.
  \yes{C} This is the answer: always taking the single highest-scoring molecule \emph{increases} exploitation and worsens rediscovery, so it is \emph{not} a remedy.
\end{answer}

\questiontitle{Autoregressive molecular reinforcement learning}
When molecule construction is performed autoregressively and treated as a
reinforcement-learning problem, which assignment of the RL roles is correct?
\begin{itemize}[leftmargin=1.8em]
  \choice The partial molecule is the state and the next building step is the action, but the scoring function is the policy and the generator gives the reward.
  \choice The finished molecule is the state, the scoring function is the action, and each partial molecule gets its own final reward.
  \choice The partial molecule is the state, the next building step is the action, the generator is the policy, and the finished molecule earns the reward.
\end{itemize}
\closequestion
\begin{answer}{C}
  \no{A} The \emph{generator} is the policy, not the scoring function, and the generator does not ``supply the reward''.
  \no{B} The completed molecule is not the state, and partial molecules do not each get an independent final reward.
  \yes{C} Partial molecule $=$ state, next construction step $=$ action, generator $=$ policy, completed molecule earns the reward.
\end{answer}

\questiontitle{Diffusion vs.\ one-shot latent decoding}
What sets diffusion-based molecular generation apart from decoding a molecule from a
single latent vector in one pass?
\begin{itemize}[leftmargin=1.8em]
  \choice It maps a single latent vector straight to a finished molecule in one decoder pass.
  \choice It begins from noise and repeatedly applies a learned denoising process.
  \choice It builds the molecule step by step, predicting the next token or graph action from the partial molecule.
\end{itemize}
\closequestion
\begin{answer}{B}
  \no{A} Mapping one latent vector to a full molecule in a single pass is exactly one-shot decoding.
  \yes{B} Diffusion starts from noise and \emph{iteratively} denoises toward a molecule.
  \no{C} Step-by-step next-token / graph-action generation is autoregressive, not diffusion.
\end{answer}

\questiontitle{Conditional generation}
Trained on molecule--logP pairs, a model takes a target logP value as input and emits
molecules in a single forward pass, never querying a logP scoring function during
generation. Which generative paradigm is this?
\begin{itemize}[leftmargin=1.8em]
  \choice Conditional generation.
  \choice Goal-directed generation.
  \choice Distribution learning followed by unconditional sampling.
\end{itemize}
\closequestion
\begin{answer}{A}
  \yes{A} The property is given as a \emph{condition}, the model maps it to molecules in one pass, and no scoring function is queried --- conditional generation.
  \no{B} Goal-directed generation optimises against a scoring function during generation, which is explicitly excluded here.
  \no{C} Unconditional sampling would ignore the desired logP input.
\end{answer}

\questiontitle{Distribution learning}
In molecular generation, which option most accurately captures what distribution learning is?
\begin{itemize}[leftmargin=1.8em]
  \choice Distribution learning produces molecules conditioned on side information such as protein targets or desired property values.
  \choice Distribution learning seeks to produce molecules resembling the training set by learning its underlying distribution.
  \choice Distribution learning concentrates on producing molecules optimised for specific biological or physicochemical properties.
\end{itemize}
\closequestion
\begin{answer}{B}
  \no{A} Conditioning on auxiliary information is \emph{conditional} generation.
  \yes{B} Distribution learning models the training distribution and samples molecules resembling it.
  \no{C} Optimising for specific properties is \emph{goal-directed} generation.
\end{answer}

\questiontitle{Mode collapse}
In a generative model for molecules, what does the term mode collapse refer to?
\begin{itemize}[leftmargin=1.8em]
  \choice The model produces molecules that are valid and novel but match the training-set property distribution poorly.
  \choice The model keeps emitting the same or nearly the same molecules.
  \choice The model emits almost only molecules that already appear in the training set.
\end{itemize}
\closequestion
\begin{answer}{B}
  \no{A} Good validity/novelty but poor distribution match is a different failure mode, not mode collapse.
  \yes{B} Mode collapse: the generator keeps producing the same or very similar molecules (it covers only a few ``modes'').
  \no{C} Re-emitting training molecules is memorisation, not mode collapse.
\end{answer}

\questiontitle{Reaction prediction and retrosynthesis}
Which of the options correctly pairs each task with its corresponding inputs and outputs?
\begin{itemize}[leftmargin=1.8em]
  \choice Both forward prediction and single-step retrosynthesis take only the reaction conditions and predict the yield.
  \choice Forward prediction maps reactants to products, while single-step retrosynthesis maps a target product to candidate precursors.
  \choice Forward prediction maps products to reactants, while single-step retrosynthesis predicts the major product from given reactants.
\end{itemize}
\closequestion
\begin{answer}{B}
  \no{A} Neither task takes only conditions and predicts yield.
  \yes{B} Forward prediction: reactants $\to$ products; single-step retrosynthesis: product $\to$ precursors.
  \no{C} The two directions are swapped.
\end{answer}

\questiontitle{Reaction templates}
What role is a reaction template intended to serve?
\begin{itemize}[leftmargin=1.8em]
  \choice To give a general symbolic rule stating the required substructures and how reactants turn into products.
  \choice To capture reaction conditions such as solvent and temperature, leaving the structural change to a neural model.
  \choice To attach a reaction-class label grouping similar reactions, without saying how atoms or bonds change.
\end{itemize}
\closequestion
\begin{answer}{A}
  \yes{A} A reaction template is a generalised symbolic rule: required substructures plus how reactants become products.
  \no{B} Encoding solvent/temperature conditions is not what a template captures.
  \no{C} Assigning a reaction \emph{class} without the atom/bond transformation is classification, not a template.
\end{answer}

\questiontitle{Easiest reaction task}
Among the tasks listed below, which one is typically the least difficult?
\begin{itemize}[leftmargin=1.8em]
  \choice Single-step retrosynthesis.
  \choice Multi-step retrosynthesis planning.
  \choice Forward reaction prediction.
\end{itemize}
\closequestion
\begin{answer}{C}
  \no{A} Single-step retrosynthesis is one-to-many (many valid precursor sets), so it is harder.
  \no{B} Multi-step planning searches over many steps --- the hardest of the three.
  \yes{C} Forward reaction prediction is the most constrained (reactants $\to$ major product) and generally the easiest.
\end{answer}

\questiontitle{Template-free translation methods}
In what terms do template-free translation methods formulate reaction prediction or
retrosynthesis?
\begin{itemize}[leftmargin=1.8em]
  \choice As looking up the most compatible stored reaction pattern in an associative memory.
  \choice As sorting the input into a learned reaction category and then applying a symbolic transformation.
  \choice As a direct sequence-to-sequence translation between molecular reaction strings.
\end{itemize}
\closequestion
\begin{answer}{C}
  \no{A} Retrieving stored patterns from memory is a template/retrieval approach.
  \no{B} Classifying into a learned reaction class then transforming is template-based.
  \yes{C} Template-free methods translate reaction strings directly, as sequence-to-sequence (e.g.\ the Molecular Transformer).
\end{answer}

\questiontitle{Deep networks inside MCTS for retrosynthesis}
What makes deep neural networks helpful when they are embedded within Monte Carlo
Tree Search for retrosynthesis?
\begin{itemize}[leftmargin=1.8em]
  \choice They take over the search by scoring only the first disconnection and then running that route greedily.
  \choice They can estimate the value of unvisited states and propose promising actions when exhaustive search is impractical.
  \choice They mostly balance chemical equations at each node, while the search alone fixes value and policy without learned guidance.
\end{itemize}
\closequestion
\begin{answer}{B}
  \no{A} They do not collapse the search to a single greedy first disconnection.
  \yes{B} Learned networks estimate the value of unvisited states and propose promising actions, guiding the search where exhaustive enumeration is infeasible.
  \no{C} They provide the value/policy guidance --- the search is not left unguided.
\end{answer}

\questiontitle{Pocket-conditioned 3D ligand generation}
Which option most accurately describes pocket-conditioned 3D ligand generation?
\begin{itemize}[leftmargin=1.8em]
  \choice Ligands are built without protein context and only docked afterwards; the docking score does not steer the generator.
  \choice The protein pocket supplies geometric context for building or denoising a ligand, usually followed by validity, docking, property, and synthesis checks.
  \choice A fixed ligand graph is supplied and only different conformations of that same molecule are sampled, ignoring the protein.
\end{itemize}
\closequestion
\begin{answer}{B}
  \no{A} Generating without protein context and only docking afterwards is \emph{not} pocket-conditioned --- the pocket never conditions the generator.
  \yes{B} The binding pocket gives geometric context to generate/denoise the ligand, typically followed by validity, docking, property, and synthesis checks.
  \no{C} Sampling conformers of a fixed ligand without protein information is not pocket-conditioned generation.
\end{answer}


\end{document}
